% Resumen y palabras clave; Abstract y keywords (consigna 2.2, ítems 3 y 4). Se muestran en el recuadro de la
% primera página (00_portada.tex), como en las revistas académicas.
% Resumen: texto del estudiante (versión revisada del 25-nov-2026 y reescritura del 26-nov-2026).
% Abstract: traducción al inglés del resumen final, asistida por IA (26-nov-2026).
\begin{resumen}{Resumen}
Este estudio caracteriza la oferta de financiamiento participativo financiero (FPF) en el Perú entre 2022 y 2026 y la compara con las condiciones del crédito tradicional. Para ello se armó un panel reproducible de \cifra{filasPanel} observaciones plataforma-mes, con datos de las interfaces de programación del Banco Central de Reserva del Perú y de la Superintendencia del Mercado de Valores (SMV) y del rastreo del portal de proyectos de la SMV. Todas las variables usadas en las pruebas econométricas tienen como mínimo 1\,000 observaciones o \cifra{muSeriesObs} periodos mensuales. De las \cifra{plataformasAutorizadas} sociedades administradoras, solo \cifra{plataformasConFinanciamiento} financiaron proyectos, y una reunió el \cifra{participacionLider}~\% del monto, que subió de S/~\cifra{monto2022} millones en 2022 a S/~\cifra{monto2025} millones en 2025. Con \cifra{muSeriesObs} meses de datos, las pruebas de Engle-Granger y de Johansen no encuentran cointegración entre el monto del FPF y las tasas del crédito tradicional, y en el corto plazo esas tasas tampoco resultan significativas en conjunto ($p\cifra{cpConWaldPRel}$). Después de la modificación del reglamento de 2025 no se observa un efecto positivo, porque el monto ya venía creciendo desde antes. Un proyecto empresarial tiene mayor probabilidad de ser financiado (razón de momios de \cifra{logUnoOr}), aunque esa ventaja se reduce a \cifra{logTresOr} cuando se controla por la plataforma y el año. Al controlar el tamaño, la rentabilidad de las plataformas no se diferencia de manera significativa de la de otros intermediarios ($p\cifra{empRoaPRel}$). Como la SMV no publica la tasa ni el plazo de las operaciones, el precio del FPF no se puede comparar con el del crédito tradicional.
\palabrasclave{Palabras clave}{crowdfunding, financiamiento participativo financiero, fintech, cointegración, Perú}
\end{resumen}%
%
\begin{resumen}{Abstract}
This study characterizes the supply of lending-based crowdfunding in Peru between 2022 and 2026 and compares it with the terms of traditional credit. A reproducible panel of \cifra{filasPanel} platform-month observations was built with data from the application programming interfaces of the Central Reserve Bank of Peru and the Superintendency of the Securities Market (SMV), and from web scraping of the SMV project portal. All the variables used in the econometric tests have at least 1\,000 observations or \cifra{muSeriesObs} monthly periods. Only \cifra{plataformasConFinanciamiento} of the \cifra{plataformasAutorizadas} platform operators funded projects, and one of them held \cifra{participacionLider}~\% of the amount, which rose from S/~\cifra{monto2022} million in 2022 to S/~\cifra{monto2025} million in 2025. With \cifra{muSeriesObs} months of data, the Engle-Granger and Johansen tests find no cointegration between the crowdfunding amount and traditional lending rates, and in the short run those rates are not jointly significant either ($p\cifra{cpConWaldPRel}$). No positive effect is observed after the 2025 regulatory amendment, because the amount was already growing beforehand. A business project is more likely to be funded (odds ratio of \cifra{logUnoOr}), although this advantage falls to \cifra{logTresOr} when platform and year are controlled for. Once size is controlled for, platform profitability does not differ significantly from that of other intermediaries ($p\cifra{empRoaPRel}$). Because the SMV does not publish the rate or term of each loan, the price of crowdfunding cannot be compared with that of traditional credit.
\palabrasclave{Keywords}{lending-based crowdfunding, fintech, cointegration, Peru}
\end{resumen}
